% \noindent 
% Questo modello Agent-Based (ABM) simula le dinamiche di trading di un asset (ispirato al Bitcoin) utilizzando uno stormo di agenti divisi in tre categorie (Breed). Il modello evolve in due fasi: una \textbf{Fase 1 (Esogena)} guidata dai dati reali per calibrare il mercato, e una \textbf{Fase 2 (Endogena)} in cui il prezzo e i volumi sono determinati interamente dalle interazioni degli agenti.

% \section{Topologia e Comportamento degli Agenti}
% Gli agenti si muovono in uno spazio 2D che rappresenta il loro ``stato mentale'' rispetto al mercato:
% \begin{itemize}
%     \item \textbf{Asse Y (Sentiment):} Determina la direzione del trade. Se $Y \ge 0$, l'agente è propenso all'acquisto (Buy). Se $Y < 0$, è propenso alla vendita (Sell).
%     \item \textbf{Asse X (Convinzione):} Determina la probabilità effettiva di eseguire il trade. La probabilità $p$ cresce man mano che l'agente si sposta verso destra:
% \end{itemize}

% \begin{equation}
%     p = \frac{X - X_{min}}{X_{max} - X_{min}}
% \end{equation}

% \noindent Gli agenti sono divisi in tre tipologie di comportamento rispetto ai rendimenti (Return) del mercato:
% \begin{enumerate}
%     \item \textbf{Specie A (Gialli):} Trend-follower (seguono il mercato).
%     \item \textbf{Specie B (Rossi):} Contrarian (vanno contro il mercato).
%     \item \textbf{Specie C (Ciano):} Rumore di fondo (non influenzati direttamente dal rendimento per la direzione Y).
% \end{enumerate}

% \section{Dinamica del Movimento (Corrente e Hype)}
% Gli agenti vengono spinti da una ``corrente'' macroeconomica. La componente orizzontale (volatilità) spinge tutti verso una maggiore probabilità di trade:

% \begin{equation}
%     C_x = k \cdot \text{Volatilità}
% \end{equation}

% \noindent La componente verticale dipende dalla specie e dal rendimento (Return):
% \begin{align}
%     C_y^{(A)} &= h \cdot \text{Return} \\
%     C_y^{(B)} &= -h \cdot \text{Return} \\
%     C_y^{(C)} &= 0
% \end{align}

% \subsection*{Hype e Stubbornness (Testardaggine)}
% L'Hype (stress di mercato) è definito come la somma del rendimento assoluto e della volatilità:
% \begin{equation}
%     H = |\text{Return}| + \text{Volatilità}
% \end{equation}

% \noindent Quando l'Hype sale, gli agenti diventano meno ``testardi'' ($S_{dyn}$) e più inclini a seguire il gregge (flocking):
% \begin{equation}
%     S_{dyn} = \max\left(S_{min}, \min\left(S_{base}, S_{base} - (K_{hype} \cdot H)\right)\right)
% \end{equation}
% Dove $K_{hype}$ è la sensibilità all'hype impostata dall'utente.

% \section{Transizione alla Fase 2 e Prezzo Endogeno}
% Il modello calcola la correlazione di Pearson $r$ tra il volume generato dagli agenti e il volume reale. Se $r \ge \text{Threshold}$ per più di 5 tick consecutivi, il modello ``sgancia'' i dati reali e inizia a generare il prezzo endogenamente.

% \vspace{0.3cm}
% \noindent Il nuovo prezzo viene calcolato in base alla \textbf{Pressione di Mercato} ($P_{m}$), definita come la media delle coordinate Y degli agenti rispetto al massimo possibile:
% \begin{equation}
%     P_{m} = \frac{\bar{Y}}{Y_{max}}
% \end{equation}

% \noindent Il prezzo endogeno al tempo $t+1$ si aggiorna usando il parametro \texttt{price-sensitivity} ($S_p$):
% \begin{equation}
%     \text{Price}_{t+1} = \max\left(1, \text{Price}_t \cdot (1 + S_p \cdot P_{m})\right)
% \end{equation}

% \section{Il Doppio Motore di Isteresi (Memoria di Mercato)}
% L'isteresi simula il fatto che i volumi in dollari non dipendono solo dall'azione istantanea, ma mantengono ``memoria'' dei picchi storici raggiunti (All-Time High). Nel modello agiscono due forze di isteresi moltiplicative.

% \subsection{Isteresi del Prezzo ($\beta_{vol}$)}
% Misura quanto il mercato ricorda i volumi basandosi sul crollo del prezzo rispetto al suo picco. Siano:
% \begin{equation}
%     R_{p}^{(current)} = \frac{\text{Price}_{corrente}}{\text{Price}_{iniziale}}, \quad R_{p}^{(ath)} = \frac{\text{Price}_{ath}}{\text{Price}_{iniziale}}
% \end{equation}
% Il fattore di isteresi del prezzo $F_p$ è:
% \begin{equation}
%     F_p = \max \left( 0.01, R_{p}^{(current)} + \beta_{vol} \left( R_{p}^{(ath)} - R_{p}^{(current)} \right) \right)
% \end{equation}

% \subsection{Isteresi del Volume ($\gamma_{vol}$)}
% Misura quanto il mercato mantiene i volumi basandosi puramente sulle masse scambiate in passato (memoria dei flussi). Siano:
% \begin{equation}
%     R_{v}^{(current)} = \frac{\text{Volume}_{corrente}}{\text{Volume}_{iniziale}}, \quad R_{v}^{(ath)} = \frac{\text{Volume}_{ath}}{\text{Volume}_{iniziale}}
% \end{equation}
% Il fattore di isteresi del volume $F_v$ è:
% \begin{equation}
%     F_v = \max \left( 0.01, R_{v}^{(current)} + \gamma_{vol} \left( R_{v}^{(ath)} - R_{v}^{(current)} \right) \right)
% \end{equation}

% \subsection{Volume Finale in Dollari}
% Il volume totale scambiato in USD dal modello, che tiene conto del volume grezzo scambiato dagli agenti ($Vol_{raw}$) moltiplicato per il prezzo attuale e ``gonfiato'' dalla memoria del mercato, è infine:
% \begin{equation}
%     \text{Vol}_{USD} = \text{Vol}_{raw} \cdot \text{Price}_{corrente} \cdot F_p \cdot F_v
% \end{equation}

\noindent
Il progetto sviluppa un modello \textit{agent-based} per simulare la formazione
di dinamiche di trading su un asset digitale ispirato a Bitcoin. L'obiettivo non
è costruire un modello predittivo del prezzo reale, ma realizzare un ambiente
computazionale in cui il comportamento collettivo degli agenti possa generare
grandezze aggregate confrontabili con quelle osservate nei dati di mercato:
prezzo, volume, volatilità, correlazioni e grado di coordinamento della
popolazione.

Il modello combina due livelli descrittivi. Da un lato è presente una componente
finanziaria, costituita da prezzo, rendimento, volatilità, volume di scambio e
portafoglio degli agenti. Dall'altro lato è presente una componente dinamica di
tipo \textit{flocking}, in cui gli agenti si muovono in uno spazio bidimensionale,
interagiscono con i propri vicini e modificano il proprio comportamento in
funzione dello stato del mercato.

Una caratteristica centrale del modello è la presenza di due fasi. Nella
\textbf{Fase 1}, detta esogena, il prezzo del modello viene agganciato al prezzo
reale o sintetico fornito dalla serie storica. In questa fase il sistema usa il
dato esterno come riferimento iniziale. Nella \textbf{Fase 2}, detta endogena, il
prezzo non viene più copiato dai dati, ma viene prodotto dalla dinamica interna
degli agenti. Il passaggio tra le due fasi avviene automaticamente quando il
volume generato dal modello risulta sufficientemente correlato con il volume
reale.

\section{Struttura generale del modello}

Il modello considera tre popolazioni di agenti, indicate come \texttt{agents-a},
\texttt{agents-b} e \texttt{agents-c}. Le tre specie non rappresentano singoli
investitori reali, ma categorie comportamentali astratte. Esse permettono di
studiare come gruppi con risposte diverse allo stesso segnale di mercato possano
produrre dinamiche aggregate differenti.

Ogni agente possiede alcune variabili individuali fondamentali:
\begin{itemize}
    \item una posizione nello spazio bidimensionale, data dalle coordinate
    \texttt{xcor} e \texttt{ycor};
    \item una velocità, descritta dalle componenti \texttt{velocity-x} e
    \texttt{velocity-y};
    \item un insieme di vicini o \texttt{flockmates}, che definisce la rete di
    riferimento per il meccanismo di allineamento;
    \item un portafoglio finanziario, composto da una quantità di Bitcoin
    detenuta, \texttt{bitcoin-held}, e da una disponibilità liquida,
    \texttt{cash-balance}.
\end{itemize}

Durante il \texttt{setup}, gli agenti vengono creati in una regione circolare
centrale, ricevono una direzione iniziale casuale e vengono inizializzati con una
dotazione casuale di Bitcoin e liquidità. Inoltre, per ogni specie viene costruita
una rete fissa di interazione: ogni agente seleziona un certo numero di
\texttt{flockmates} appartenenti alla stessa specie. In questo modo il modello
separa la vicinanza sociale, rappresentata dalla rete, dalla vicinanza spaziale,
rappresentata dalla posizione nello spazio.

\section{Interpretazione dello spazio comportamentale}

Lo spazio bidimensionale in cui si muovono gli agenti non rappresenta uno spazio
fisico, ma uno spazio comportamentale. Le coordinate dell'agente sintetizzano il
suo stato rispetto al mercato.

L'asse verticale \(y\) rappresenta la direzione dell'azione di trading:
\begin{itemize}
    \item se \(y \geq 0\), l'agente si trova nella regione di acquisto;
    \item se \(y < 0\), l'agente si trova nella regione di vendita.
\end{itemize}

L'asse orizzontale \(x\) rappresenta invece la propensione operativa dell'agente,
cioè la probabilità che esso esegua effettivamente una compravendita. Tale
probabilità è calcolata normalizzando la posizione orizzontale dell'agente
rispetto alla larghezza del mondo:

\begin{equation}
    p_i =
    \frac{x_i - x_{\min}}{x_{\max} - x_{\min}}.
\end{equation}

Nel codice il valore viene ulteriormente limitato all'intervallo \([0,1]\), in
modo da ottenere sempre una probabilità valida:

\begin{equation}
    p_i =
    \min\left[
    1,
    \max\left(
    0,
    \frac{x_i - x_{\min}}{x_{\max} - x_{\min}}
    \right)
    \right].
\end{equation}

Questa scelta permette di attribuire un significato comportamentale semplice alle
coordinate: la posizione verticale stabilisce se l'agente tende a comprare o a
vendere, mentre la posizione orizzontale stabilisce con quale probabilità
partecipa al mercato.

\section{Specie di agenti e risposta al rendimento}

Le tre specie di agenti differiscono per il modo in cui rispondono al rendimento
del mercato. Il rendimento attivo può essere quello reale, nella Fase 1, oppure
quello endogeno, nella Fase 2. Indichiamo tale rendimento con \(r(t)\).

La componente verticale della corrente è definita in modo diverso per le tre
specie:

\begin{align}
    J_y^{(A)}(t) &= h r(t), \\
    J_y^{(B)}(t) &= -h r(t), \\
    J_y^{(C)}(t) &= 0.
\end{align}

La specie A può essere interpretata come una popolazione \textit{trend-following}:
quando il rendimento è positivo, essa viene spinta verso la regione di acquisto;
quando il rendimento è negativo, viene spinta verso la regione di vendita.

La specie B ha una risposta opposta e può essere interpretata come una popolazione
\textit{contrarian}: tende a muoversi contro il segnale di rendimento. La specie C
non riceve una spinta verticale diretta dal rendimento e funziona come componente
neutrale o rumore di fondo.

Questa distinzione consente di osservare come l'interazione tra gruppi con
comportamenti diversi influenzi la dinamica complessiva del mercato simulato.

\section{Corrente di mercato e dinamica del movimento}

Il movimento degli agenti è influenzato da una corrente di mercato. La corrente
ha una componente orizzontale e una componente verticale. La componente
orizzontale dipende dalla volatilità attiva:

\begin{equation}
    J_x(t) = k \sigma(t),
\end{equation}

dove \(k\) è un parametro di intensità e \(\sigma(t)\) è la volatilità. Nella
Fase 1 viene usata la volatilità reale, mentre nella Fase 2 viene usata la
volatilità endogena del modello.

La funzione della componente orizzontale è spingere gli agenti lungo l'asse della
propensione operativa. Quando la volatilità è più alta, il mercato è più
instabile e gli agenti tendono a muoversi verso regioni in cui la probabilità di
trading è maggiore.

Il movimento è descritto mediante una velocità vettoriale. A ogni passo, la
corrente modifica le componenti della velocità:

\begin{align}
    v_x(t+1) &= v_x(t) + J_x(t), \\
    v_y(t+1) &= v_y(t) + J_y(t).
\end{align}

La nuova posizione prevista è poi calcolata come

\begin{align}
    x_i(t+1) &= x_i(t) + v_{x,i}(t+1), \\
    y_i(t+1) &= y_i(t) + v_{y,i}(t+1).
\end{align}

Il modello utilizza condizioni al bordo riflettenti. Se un agente raggiunge un
bordo del mondo, la componente corrispondente della velocità viene invertita.
Questa scelta mantiene gli agenti confinati nello spazio comportamentale ed evita
che escano dal dominio della simulazione.

Infine, viene introdotta una perturbazione casuale della direzione di moto. La
velocità viene ruotata di un angolo casuale compreso tra
\(-\texttt{temperature}\) e \(+\texttt{temperature}\). Il parametro
\texttt{temperature} controlla quindi il grado di rumore della dinamica: valori
bassi producono traiettorie più regolari, mentre valori alti rendono il moto più
disordinato.

\section{Flocking, hype e stubbornness}

\subsection{Condizioni iniziali}
La scelta di inizializzare gli agenti in una regione circolare centrale non deve
essere interpretata come una descrizione realistica dello stato iniziale del
mercato, ma come una condizione iniziale neutra. Essa rappresenta una popolazione
inizialmente poco polarizzata, nella quale gli agenti non sono ancora fortemente
orientati né verso l'acquisto né verso la vendita.

Questa scelta è utile per studiare se le dinamiche collettive osservate nel
modello emergano dalle regole di interazione, dalla corrente di mercato e dal
meccanismo di flocking, piuttosto che essere imposte artificialmente dalla
condizione iniziale.

Tuttavia, la condizione iniziale non è unica. Per analizzare la robustezza del
modello è opportuno confrontare diverse inizializzazioni: distribuzione centrale,
distribuzione uniforme, distribuzioni sbilanciate verso l'acquisto o la vendita,
e condizioni differenziate per specie. In questo modo è possibile distinguere gli
effetti realmente generati dalla dinamica del modello da quelli dipendenti dalla
scelta iniziale.\\


Gli agenti non si muovono in modo completamente indipendente. Ogni agente tende
ad allineare la propria direzione con quella dei propri \texttt{flockmates}. Il
nuovo orientamento viene calcolato come combinazione pesata tra la direzione
individuale dell'agente e la direzione media dei suoi vicini.

Sia \(\theta_i\) la direzione dell'agente e siano

\begin{equation}
    d_{x,i} = \cos(\theta_i),
    \qquad
    d_{y,i} = \sin(\theta_i)
\end{equation}

le componenti della sua direzione. Indicando con
\(\bar{d}_{x,i}\) e \(\bar{d}_{y,i}\) le componenti medie dei suoi
\texttt{flockmates}, la direzione aggiornata è costruita come

\begin{align}
    d_{x,i}^{new}
    &=
    s(t)d_{x,i}
    +
    \left[1-s(t)\right]\bar{d}_{x,i}, \\
    d_{y,i}^{new}
    &=
    s(t)d_{y,i}
    +
    \left[1-s(t)\right]\bar{d}_{y,i}.
\end{align}

Il parametro \(s(t)\) rappresenta la \textit{stubbornness}, cioè la tendenza
dell'agente a mantenere la propria direzione precedente. Se \(s(t)\) è alto,
l'agente è poco influenzato dal gruppo; se \(s(t)\) è basso, l'agente tende ad
allinearsi maggiormente ai propri vicini.

Nel modello la stubbornness non è costante, ma dipende dall'hype di mercato. Lo
hype è definito come

\begin{equation}
    H(t)=|r_m(t)|+\sigma_m(t),
\end{equation}

dove \(r_m(t)\) è il rendimento endogeno del modello e \(\sigma_m(t)\) è la
volatilità endogena. Quando il mercato è più instabile, cioè quando rendimento e
volatilità sono più elevati, gli agenti diventano meno testardi e più sensibili
al comportamento collettivo.

La stubbornness dinamica è calcolata come

\begin{equation}
    s(t)
    =
    \max
    \left[
    s_{\min},
    \min
    \left(
    s_{\mathrm{base}},
    s_{\mathrm{base}} -
    K_{\mathrm{hype}}H(t)
    \right)
    \right].
\end{equation}

Il parametro \(K_{\mathrm{hype}}\) controlla quanto rapidamente la stubbornness
diminuisce all'aumentare dello stress di mercato.

\section{Trading e portafoglio degli agenti}

Ogni agente possiede un portafoglio formato da liquidità e Bitcoin. La procedura
di trading utilizza il prezzo endogeno corrente del modello, indicato con
\(P_m(t)\). Se il prezzo è positivo, l'agente può comprare o vendere a seconda
della propria posizione verticale.

Se \(y_i \geq 0\), l'agente si trova nella regione di acquisto. La quantità
massima acquistabile è

\begin{equation}
    q_i^{buy}
    =
    \min
    \left(
    \texttt{trade-size},
    \left\lfloor
    \frac{C_i}{P_m(t)}
    \right\rfloor
    \right),
\end{equation}

dove \(C_i\) è la liquidità disponibile dell'agente. Questa formula impedisce
all'agente di comprare più Bitcoin di quanto possa permettersi.

Se \(y_i < 0\), l'agente si trova nella regione di vendita. La quantità massima
vendibile è

\begin{equation}
    q_i^{sell}
    =
    \min
    \left(
    \texttt{trade-size},
    B_i
    \right),
\end{equation}

dove \(B_i\) è la quantità di Bitcoin detenuta dall'agente. In questo modo non è
possibile vendere Bitcoin non posseduti.

L'operazione avviene solo se un numero casuale uniforme tra 0 e 1 è minore della
probabilità di scambio \(p_i\). Ogni acquisto o vendita riuscita contribuisce al
volume di trading del modello, misurato inizialmente in Bitcoin.

\section{Dati di mercato e scenari}

Il modello può essere eseguito utilizzando diverse sorgenti di dati. La procedura
\texttt{scarica-dati} seleziona lo scenario in base alla variabile
\texttt{scenario}. Le possibilità previste sono:

\begin{itemize}
    \item caricamento di una serie storica da file CSV;
    \item utilizzo di un prezzo live tramite API;
    \item generazione di scenari sintetici di mercato.
\end{itemize}

Gli scenari sintetici includono:
\begin{itemize}
    \item \textit{bull market}, con drift positivo;
    \item \textit{bear market}, con drift negativo;
    \item mercato laterale, con drift nullo;
    \item crash improvviso;
    \item bolla speculativa seguita da collasso.
\end{itemize}

Questa struttura rende il modello flessibile: è possibile confrontare la risposta
degli agenti sia a dati esterni sia a condizioni artificiali controllate.

\section{Fase 1: dinamica esogena}

Nella Fase 1 il modello è guidato dai dati esterni. Il prezzo endogeno viene
posto uguale al prezzo reale o sintetico letto dalla serie storica:

\begin{equation}
    P_m(t) = P_{BTC}(t).
\end{equation}

Anche rendimento e volatilità del modello vengono copiati dai dati esterni:

\begin{align}
    r_m(t) &= r_{BTC}(t), \\
    \sigma_m(t) &= \sigma_{BTC}(t).
\end{align}

Questa fase ha la funzione di inizializzare il sistema e permettere agli agenti
di produrre volumi confrontabili con quelli della serie di riferimento. Il
modello resta in questa fase finché la correlazione media tra volume simulato e
volume reale non supera stabilmente una soglia prefissata.

\section{Correlazione dei volumi e passaggio alla Fase 2}

Il modello confronta il volume reale con il volume generato dagli agenti mediante
la correlazione di Pearson. Date due serie \(X\) e \(Y\), la correlazione è

\begin{equation}
    r =
    \frac{
    \sum_i (x_i-\bar{x})(y_i-\bar{y})
    }{
    \sqrt{\sum_i (x_i-\bar{x})^2}
    \sqrt{\sum_i (y_i-\bar{y})^2}
    }.
\end{equation}

Nel caso dei volumi, \(X\) rappresenta lo storico recente dei volumi simulati,
mentre \(Y\) rappresenta lo storico recente dei volumi reali. Il calcolo viene
effettuato su una finestra mobile di ampiezza \texttt{window-size}, così da
misurare la somiglianza recente tra le due serie.

Il passaggio alla Fase 2 avviene quando la correlazione media dei volumi supera
la soglia \texttt{pearson-threshold} per più di cinque tick consecutivi:

\begin{equation}
    \overline{r}_V \geq \texttt{pearson-threshold}.
\end{equation}

Quando questa condizione viene soddisfatta, la variabile
\texttt{current-disabled?} viene impostata a \texttt{true}. Da quel momento il
modello smette di copiare il prezzo esterno e inizia a generare il prezzo in modo
endogeno.

\section{Fase 2: prezzo endogeno}

Nella Fase 2 il prezzo viene determinato dalla posizione media verticale degli
agenti. La coordinata \(y\) è interpretata come indicatore della pressione di
acquisto o vendita. La pressione di mercato è definita come

\begin{equation}
    M(t)=\frac{\bar{y}(t)}{y_{\max}},
\end{equation}

dove \(\bar{y}(t)\) è la media delle coordinate verticali di tutti gli agenti e
\(y_{\max}\) è il massimo valore verticale del mondo NetLogo.

Se \(M(t)>0\), la popolazione è mediamente orientata verso l'acquisto e il prezzo
subisce una pressione rialzista. Se \(M(t)<0\), la popolazione è mediamente
orientata verso la vendita e il prezzo subisce una pressione ribassista.

Il prezzo endogeno viene aggiornato secondo

\begin{equation}
    P_m(t+1)
    =
    \max
    \left[
    1,
    P_m(t)
    \left(
    1+
    S_p M(t)
    \right)
    \right],
\end{equation}

dove \(S_p=\texttt{price-sensitivity}\) controlla la sensibilità del prezzo alla
pressione di mercato. Il limite inferiore pari a 1 evita che il prezzo diventi
nullo o negativo.

Il rendimento endogeno viene poi calcolato come

\begin{equation}
    r_m(t)
    =
    \frac{P_m(t)-P_m(t-1)}{P_m(t-1)}.
\end{equation}

La volatilità endogena è approssimata dal valore assoluto del rendimento, con una
soglia minima:

\begin{equation}
    \sigma_m(t)
    =
    \max
    \left(
    0.005,
    |r_m(t)|
    \right).
\end{equation}

La soglia minima impedisce alla volatilità di annullarsi completamente e mantiene
attivi i meccanismi dipendenti da essa.

\section{Volume simulato e isteresi}

Il volume grezzo del modello, \(\text{Vol}_{raw}\), è dato dalla quantità totale
di Bitcoin scambiata dagli agenti nel tick corrente. Per confrontarlo con un
volume reale espresso in dollari, esso viene moltiplicato per il prezzo endogeno:

\begin{equation}
    \text{Vol}_{raw}^{USD}
    =
    \text{Vol}_{raw} \cdot P_m(t).
\end{equation}

Nel modello viene introdotto un fattore di memoria aggregata del mercato. Tale
fattore è ispirato al concetto di isteresi, intesa come dipendenza dello stato
attuale di un sistema non solo dal valore corrente delle variabili rilevanti, ma
anche dal percorso con cui tale stato è stato raggiunto.

Nel caso considerato, il volume simulato non dipende soltanto dal prezzo corrente,
ma anche dal massimo prezzo raggiunto nella finestra di simulazione. Due stati con
lo stesso prezzo corrente possono quindi produrre livelli diversi di attività di
mercato se sono stati preceduti da storie differenti. In questo senso il termine
può essere interpretato come una forma fenomenologica di path dependence.

È importante sottolineare che tale memoria non rappresenta necessariamente una
memoria esplicita dei singoli agenti. Nel modello attuale essa è introdotta come
variabile macroscopica aggregata, utile a rappresentare la persistenza
dell'attenzione e dell'attività di mercato dopo fasi di forte crescita del prezzo.
Siano

\begin{equation}
    R_{\mathrm{current}}
    =
    \frac{P_m(t)}{P_{\mathrm{initial}}},
    \qquad
    R_{\mathrm{ATH}}
    =
    \frac{P_{\mathrm{ATH}}}{P_{\mathrm{initial}}},
\end{equation}

dove \(P_{\mathrm{initial}}\) è il prezzo iniziale e \(P_{\mathrm{ATH}}\) è il
massimo prezzo raggiunto nella finestra di simulazione. Il fattore di isteresi è

\begin{equation}
    H_V(t)
    =
    \max
    \left[
    0.01,
    R_{\mathrm{current}}
    +
    \beta_{\mathrm{vol}}
    \left(
    R_{\mathrm{ATH}} -
    R_{\mathrm{current}}
    \right)
    \right].
\end{equation}

Il parametro \(\beta_{\mathrm{vol}}\) controlla il peso della memoria del massimo
storico. Se \(\beta_{\mathrm{vol}}\) è vicino a zero, il fattore dipende quasi
solo dal prezzo corrente. Se invece è più alto, il modello mantiene una memoria
più marcata dei picchi passati.

Il volume finale in dollari prodotto dal modello è quindi

\begin{equation}
    \text{Vol}_{model}^{USD}
    =
    \text{Vol}_{raw}
    \cdot
    P_m(t)
    \cdot
    H_V(t).
\end{equation}

\section{Metriche collettive e stabilità}

Oltre a prezzo e volume, il modello calcola alcune metriche utili per analizzare
la struttura collettiva degli agenti.

La prima è il parametro d'ordine, che misura il grado di allineamento della
popolazione. Per un insieme di \(N\) agenti, esso è definito come

\begin{equation}
    \Phi
    =
    \frac{
    \sqrt{
    \left(\sum_i dx_i\right)^2
    +
    \left(\sum_i dy_i\right)^2
    }
    }{N}.
\end{equation}

Se gli agenti sono orientati in direzioni casuali, le componenti vettoriali si
compensano e \(\Phi\) tende a valori bassi. Se invece gli agenti sono allineati,
\(\Phi\) si avvicina a 1. Il modello calcola sia il parametro d'ordine globale,
sia quello specifico per ciascuna specie.

Una seconda metrica è

\begin{equation}
    \lambda =
    \text{slope}
    \left[
    \log(\Phi(t)+0.0001)
    \right],
\end{equation}

dove la pendenza è calcolata tramite regressione lineare su una finestra mobile.
Questa grandezza fornisce una misura della tendenza recente del parametro
d'ordine: se \(\lambda>0\), l'ordine collettivo tende ad aumentare; se
\(\lambda<0\), tende a diminuire.

Il modello calcola inoltre la correlazione di fase tra le specie A e B,
utilizzando le loro posizioni medie verticali:

\begin{equation}
    r_{AB}
    =
    \mathrm{corr}
    \left(
    \bar{y}_A,
    \bar{y}_B
    \right).
\end{equation}

Questa quantità misura se le due specie si muovono nella stessa direzione lungo
l'asse comportamentale di acquisto-vendita oppure se tendono a muoversi in modo
opposto.

Infine, viene calcolato un indice basato sulla distribuzione relativa delle tre
popolazioni:

\begin{equation}
    H =
    s_A^2+s_B^2+s_C^2,
\end{equation}

dove \(s_A\), \(s_B\) e \(s_C\) sono le quote relative delle tre specie. Questo
indice è minimo quando le tre popolazioni sono bilanciate e aumenta quando una
specie domina sulle altre. Per questo motivo può essere interpretato come un
indice di concentrazione della composizione della popolazione.

\section{Sintesi del funzionamento del modello}

Il funzionamento complessivo del modello può essere riassunto come segue. In una
prima fase, gli agenti vengono inizializzati con portafogli eterogenei e si
muovono in uno spazio comportamentale influenzato da segnali di mercato esterni.
La loro posizione determina la probabilità e la direzione delle operazioni di
trading. Le operazioni aggregate producono un volume simulato, che viene
confrontato con il volume reale tramite correlazione di Pearson.

Quando la correlazione dei volumi supera stabilmente la soglia impostata, il
modello passa alla fase endogena. Da quel momento il prezzo non è più imposto
dall'esterno, ma emerge dalla distribuzione degli agenti nello spazio. Il mercato
simulato diventa quindi il risultato dell'interazione tra movimento, flocking,
portafogli, segnali di rendimento, volatilità e memoria dei picchi passati.

Il modello costituisce quindi un laboratorio computazionale per studiare come
semplici regole individuali possano generare dinamiche aggregate complesse e
come un sistema inizialmente guidato dai dati possa evolvere verso una dinamica
autonoma.